% !TEX program = pdflatex
% ==========================================================
% CLASE BEAMER MASTER+
% Fundamentos de Matemática Básica - Agronomía
% Tema: Potencias
% Universidad Santo Tomás - Talca
% ==========================================================

\newif\ifhandout
%\handoutfalse
\handouttrue

\ifhandout
\documentclass[aspectratio=169,10pt,handout]{beamer}
\else
\documentclass[aspectratio=169,10pt]{beamer}
\fi

% ----------------------------------------------------------
% PAQUETES
% ----------------------------------------------------------
\usepackage[spanish,es-noshorthands]{babel}
\usepackage[T1]{fontenc}
\usepackage[utf8]{inputenc}
\usepackage{lmodern}
\usepackage{amsmath,amsfonts,amssymb}
\usepackage{ragged2e}
\usepackage{enumitem}
\usepackage{multicol}
\usepackage{array}
\usepackage{booktabs}
\usepackage{colortbl}
\usepackage{tikz}
\usetikzlibrary{positioning,calc,shadows.blur}
\usepackage{fontawesome5}

% ----------------------------------------------------------
% COLORES
% ----------------------------------------------------------
\definecolor{USTBlue}{RGB}{0,70,127}
\definecolor{USTBlueDark}{RGB}{0,45,85}
\definecolor{USTTeal}{RGB}{0,153,117}
\definecolor{USTGray}{RGB}{120,120,120}
\definecolor{USTSoft}{RGB}{248,249,250}
\definecolor{USTText}{RGB}{35,40,45}
\definecolor{USTGold}{RGB}{196,154,70}
\definecolor{USTLightBlue}{RGB}{227,239,248}
\definecolor{USTLightTeal}{RGB}{226,245,240}
\definecolor{USTRedSoft}{RGB}{252,235,235}
\definecolor{USTLightGold}{RGB}{247,241,224}
\definecolor{USTLine}{RGB}{223,228,233}

% ----------------------------------------------------------
% TEMA
% ----------------------------------------------------------
\usetheme{Madrid}
\usecolortheme[named=USTBlue]{structure}
\setbeamertemplate{navigation symbols}{}

% ----------------------------------------------------------
% AJUSTES VISUALES
% ----------------------------------------------------------
\setlist[itemize]{itemsep=0.45em,topsep=0.35em}
\setlist[enumerate]{itemsep=0.45em,topsep=0.35em}

\setbeamercolor{normal text}{fg=USTText,bg=white}
\setbeamercolor{title}{fg=USTBlueDark}
\setbeamercolor{subtitle}{fg=USTTeal}
\setbeamercolor{frametitle}{fg=white,bg=USTBlueDark}

\setbeamercolor{block title}{bg=USTLightBlue,fg=USTBlueDark}
\setbeamercolor{block body}{bg=USTSoft,fg=black}

\setbeamercolor{block title example}{bg=USTLightTeal,fg=USTBlueDark}
\setbeamercolor{block body example}{bg=USTSoft,fg=black}

\setbeamercolor{block title alerted}{bg=USTRedSoft,fg=black}
\setbeamercolor{block body alerted}{bg=red!2,fg=black}

\setbeamertemplate{itemize item}{\color{USTTeal}\small$\blacktriangleright$}
\setbeamertemplate{itemize subitem}{\color{USTBlueDark}\small$\bullet$}

% ----------------------------------------------------------
% FRAMETITLE
% ----------------------------------------------------------
\setbeamertemplate{frametitle}{
	\vspace*{-0.03cm}
	\begin{beamercolorbox}[wd=\paperwidth,ht=0.95cm,dp=0.2cm,leftskip=0.45cm,rightskip=0.25cm]{frametitle}
		{\large\bfseries \insertframetitle}
	\end{beamercolorbox}
}

% ----------------------------------------------------------
% PIE DE PÁGINA
% ----------------------------------------------------------
\setbeamercolor{author in head/foot}{bg=USTSoft,fg=USTGray}
\setbeamercolor{date in head/foot}{bg=USTSoft,fg=USTGray}
\setbeamercolor{title in head/foot}{bg=USTSoft,fg=USTGray}

\setbeamertemplate{footline}{%
	\leavevmode%
	\hbox{%
		\begin{beamercolorbox}[wd=.82\paperwidth,ht=2.8ex,dp=1ex,leftskip=0.4cm]{author in head/foot}%
			\scriptsize
			\textbf{Fundamentos de Matemática Básica}
			\quad|\quad
			Agronomía
			\quad|\quad
			Universidad Santo Tomás Talca
		\end{beamercolorbox}%
		\begin{beamercolorbox}[wd=.18\paperwidth,ht=2.8ex,dp=1ex,rightskip=0.35cm plus1fil]{date in head/foot}%
			\scriptsize \insertframenumber/\inserttotalframenumber
		\end{beamercolorbox}%
	}%
}

% ----------------------------------------------------------
% COMANDOS
% ----------------------------------------------------------
\newcommand{\destacar}[1]{\textcolor{USTBlueDark}{\textbf{#1}}}
\newcommand{\alerta}[1]{\textcolor{red!70!black}{\textbf{#1}}}
\newcommand{\pp}{\pause}

\newcommand{\iconoInicio}{\faPlayCircle}
\newcommand{\iconoDesarrollo}{\faBrain}
\newcommand{\iconoPractica}{\faPenSquare}
\newcommand{\iconoAgro}{\faSeedling}
\newcommand{\iconoCierre}{\faCheckCircle}
\newcommand{\iconoIdea}{\faLightbulb}
\newcommand{\iconoError}{\faExclamationTriangle}
\newcommand{\iconoTiempo}{\faClock}
\newcommand{\iconoMeta}{\faBullseye}
\newcommand{\iconoEval}{\faTasks}
\newcommand{\iconoTicket}{\faCheckCircle}
\newcommand{\iconoPuente}{\faArrowRight}

\newenvironment{metaBox}
{\begin{block}{\iconoMeta\ \ Meta de aprendizaje}}
	{\end{block}}

\newenvironment{ideaBox}
{\begin{exampleblock}{\iconoIdea\ \ Idea clave}}
	{\end{exampleblock}}

% ----------------------------------------------------------
% PORTADA MASTER+
% ----------------------------------------------------------
\newcommand{\PortadaProfesionalUST}{%
	\begin{frame}[plain]
		\begin{tikzpicture}[remember picture,overlay]
			
			\fill[white] (current page.south west) rectangle (current page.north east);
			\fill[USTBlueDark] (current page.north west) rectangle ([yshift=-1.45cm]current page.north east);
			\fill[USTTeal] ([yshift=-1.45cm]current page.north west) rectangle ([yshift=-1.62cm]current page.north east);
			\fill[USTSoft] (current page.south west) rectangle ([xshift=2.2cm]current page.north west);
			
			\fill[USTBlue!7] ([xshift=1.10cm,yshift=-2.0cm]current page.north west) circle (0.78cm);
			\fill[USTTeal!10] ([xshift=1.70cm,yshift=-4.2cm]current page.north west) circle (0.45cm);
			\fill[USTGold!14] ([xshift=-1.2cm,yshift=1cm]current page.south east) circle (1.08cm);
			
			\fill[black!8,rounded corners=8pt,blur shadow={shadow blur steps=5}]
			([xshift=-5.0cm,yshift=-3.10cm]current page.center) rectangle
			([xshift=5.0cm,yshift=2.85cm]current page.center);
			
			\fill[white,rounded corners=8pt]
			([xshift=-5.15cm,yshift=-2.95cm]current page.center) rectangle
			([xshift=4.85cm,yshift=3.0cm]current page.center);
			
			\fill[USTTeal] ([xshift=-5.15cm,yshift=-2.95cm]current page.center) rectangle
			([xshift=-4.85cm,yshift=3.0cm]current page.center);
			
			\node[anchor=north east,text=white]
			at ([xshift=-0.7cm,yshift=-0.52cm]current page.north east)
			{\small\bfseries Universidad Santo Tomás \quad Ciencias Básicas Talca};
			
			\node[align=center] at ([xshift=0.10cm,yshift=-0.02cm]current page.center){%
				\begin{minipage}{0.69\paperwidth}
					\centering
					{\Large\bfseries\color{USTText}\insertauthor\par}
					\vspace{0.75cm}
					{\fontsize{24}{28}\selectfont\bfseries\color{USTBlueDark}\inserttitle\par}
					\vspace{0.22cm}
					{\large\bfseries\color{USTTeal}\insertsubtitle\par}
					\vspace{0.55cm}
					{\normalsize\color{USTGray}\faSeedling\ Agronomía \quad \faUniversity\ Universidad Santo Tomás \quad \faChalkboardTeacher\ Clase expositiva\par}
					\vspace{0.65cm}
					{\small\color{USTGray}\insertdate\par}
				\end{minipage}
			};
			
			\draw[USTGray!35,line width=0.35pt]
			([xshift=0.9cm,yshift=0.72cm]current page.south west) --
			([xshift=-0.9cm,yshift=0.72cm]current page.south east);
			
			\node[anchor=south,text=USTGray]
			at ([yshift=0.2cm]current page.south)
			{\scriptsize \textbf{Fundamentos de Matemática Básica} \;|\; Versión MASTER+ \;|\; Potencias};
		\end{tikzpicture}
	\end{frame}
}

% ----------------------------------------------------------
% FRAME DE SECCIÓN
% ----------------------------------------------------------
\newcommand{\SectionFrame}[3]{%
	\begin{frame}[plain]
		\begin{tikzpicture}[remember picture,overlay]
			\fill[USTBlueDark] (current page.south west) rectangle (current page.north east);
			\fill[USTTeal] (current page.south west) rectangle ([yshift=0.22cm]current page.north east);
			\fill[white,opacity=0.08] ([xshift=-1cm,yshift=1cm]current page.south east) circle (1.8cm);
			\fill[white,opacity=0.06] ([xshift=1.3cm,yshift=-0.8cm]current page.north west) circle (1.25cm);
			
			\node[anchor=center,text=white] at ([yshift=0.9cm]current page.center)
			{{\fontsize{34}{36}\selectfont #1}};
			\node[anchor=center,text=white] at (current page.center)
			{{\LARGE\bfseries #2}};
			\node[anchor=center,text=white!90] at ([yshift=-1.0cm]current page.center)
			{{\large #3}};
		\end{tikzpicture}
	\end{frame}
}

% ----------------------------------------------------------
% DATOS
% ----------------------------------------------------------
\title{Fundamentos de Matemática Básica}
\subtitle{Clase MASTER+: Potencias}
\author{Prof. Luis González Alcaino}
\institute{Universidad Santo Tomás}
\date{Año académico 2026}

% ==========================================================
% DOCUMENTO
% ==========================================================
\begin{document}
	
	\PortadaProfesionalUST
	
	% ----------------------------------------------------------
	\begin{frame}{Panorama de la sesión}
		\begin{metaBox}
			Aplicar el concepto de potencia y sus propiedades en ejercicios matemáticos y situaciones de Agronomía, preparando además el paso hacia la notación científica.
		\end{metaBox}
		
		\pp
		
		\begin{block}{\faListOl\ \ Ruta de la clase}
			\begin{enumerate}
				\item Diagnóstico breve
				\item Concepto de potencia
				\item Propiedades fundamentales
				\item Errores frecuentes
				\item Práctica guiada
				\item Aplicaciones agronómicas
				\item Práctica autónoma
				\item Cierre y puente a notación científica
			\end{enumerate}
		\end{block}
	\end{frame}
	
	% ----------------------------------------------------------
	\begin{frame}{\iconoTiempo\ \ Distribución sugerida}
		\centering
		\begin{tabular}{>{\raggedright\arraybackslash}p{4.2cm} >{\centering\arraybackslash}p{2.2cm} >{\raggedright\arraybackslash}p{4.8cm}}
			\toprule
			\textbf{Momento} & \textbf{Tiempo} & \textbf{Foco} \\
			\midrule
			Diagnóstico & 8 min & Activar ideas previas \\
			Concepto y propiedades & 25 min & Comprensión conceptual \\
			Errores frecuentes & 7 min & Prevenir confusiones \\
			Práctica guiada & 15 min & Resolver paso a paso \\
			Aplicaciones agronómicas & 10 min & Transferencia profesional \\
			Práctica autónoma & 10 min & Consolidación \\
			Cierre y proyección & 5 min & Conectar con clase siguiente \\
			\bottomrule
		\end{tabular}
	\end{frame}
	
	% ----------------------------------------------------------
	\SectionFrame{\iconoInicio}{Diagnóstico inicial}{Abrimos la clase con ideas previas}
	
	% ----------------------------------------------------------
	\begin{frame}{\iconoInicio\ \ Activación}
		\begin{block}{Piensa y comenta}
			\begin{itemize}
				\item<1-> ¿Cómo escribirías \(4\cdot4\cdot4\) de forma abreviada?
				\item<2-> ¿Qué crees que significa \(10^{-3}\)?
				\item<3-> ¿Qué expresión crece más rápido: \(3n\) o \(3^n\)?
			\end{itemize}
		\end{block}
		
		\only<4>{
			\begin{ideaBox}
				El diagnóstico permite reconocer desde qué punto comienzan los estudiantes.
			\end{ideaBox}
		}
	\end{frame}
	
	% ----------------------------------------------------------
	\SectionFrame{\iconoDesarrollo}{Concepto de potencia}{Comprender antes de aplicar}
	
	% ----------------------------------------------------------
	\begin{frame}{\iconoDesarrollo\ \ ¿Qué es una potencia?}
		\centering
		\[
		a^n=\underbrace{a\cdot a\cdot a\cdots a}_{n\text{ veces}}
		\]
		
		\vspace{0.35cm}
		
		\only<2->{
			\begin{columns}[T,totalwidth=\textwidth]
				\begin{column}{0.47\textwidth}
					\begin{block}{Base}
						El número que se repite.
					\end{block}
				\end{column}
				\begin{column}{0.47\textwidth}
					\begin{exampleblock}{Exponente}
						La cantidad de repeticiones.
					\end{exampleblock}
				\end{column}
			\end{columns}
		}
		
		\only<3->{
			\begin{exampleblock}{Ejemplo}
				\[
				5^3=5\cdot5\cdot5=125
				\]
			\end{exampleblock}
		}
	\end{frame}
	
	% ----------------------------------------------------------
	\begin{frame}{\iconoDesarrollo\ \ Lectura de una potencia}
		\begin{columns}[T,totalwidth=\textwidth]
			\begin{column}{0.48\textwidth}
				\begin{block}{Se lee}
					\[
					2^4
					\]
					como “dos elevado a cuatro”.
				\end{block}
			\end{column}
			\begin{column}{0.48\textwidth}
				\begin{exampleblock}{Significa}
					\[
					2^4=2\cdot2\cdot2\cdot2=16
					\]
				\end{exampleblock}
			\end{column}
		\end{columns}
		
		\pp
		
		\begin{ideaBox}
			Una potencia no es una multiplicación cualquiera: es una multiplicación repetida de la misma base.
		\end{ideaBox}
	\end{frame}
	
	% ----------------------------------------------------------
	\SectionFrame{\iconoDesarrollo}{Propiedades}{Reglas para simplificar y calcular mejor}
	
	% ----------------------------------------------------------
	\begin{frame}{\iconoDesarrollo\ \ Propiedad 1}
		\centering
		{\Large\bfseries Producto de potencias de igual base}
		
		\vspace{0.35cm}
		
		\[
		a^m\cdot a^n=a^{m+n}
		\]
		
		\only<2->{
			\begin{exampleblock}{Ejemplo}
				\[
				2^3\cdot2^4=2^{3+4}=2^7=128
				\]
			\end{exampleblock}
		}
		
		\only<3->{
			\begin{ideaBox}
				\textbf{Clave:} si la base es la misma, se suman exponentes.
			\end{ideaBox}
		}
	\end{frame}
	
	% ----------------------------------------------------------
	\begin{frame}{\iconoDesarrollo\ \ Propiedad 2}
		\centering
		{\Large\bfseries Cociente de potencias de igual base}
		
		\vspace{0.35cm}
		
		\[
		\frac{a^m}{a^n}=a^{m-n},\qquad a\neq0
		\]
		
		\only<2->{
			\begin{exampleblock}{Ejemplo}
				\[
				\frac{5^6}{5^2}=5^{6-2}=5^4=625
				\]
			\end{exampleblock}
		}
		
		\only<3->{
			\begin{ideaBox}
				\textbf{Clave:} si la base es la misma, se restan exponentes.
			\end{ideaBox}
		}
	\end{frame}
	
	% ----------------------------------------------------------
	\begin{frame}{\iconoDesarrollo\ \ Propiedad 3}
		\centering
		{\Large\bfseries Potencia de una potencia}
		
		\vspace{0.35cm}
		
		\[
		(a^m)^n=a^{m\cdot n}
		\]
		
		\only<2->{
			\begin{exampleblock}{Ejemplo}
				\[
				(3^2)^4=3^{2\cdot4}=3^8=6561
				\]
			\end{exampleblock}
		}
		
		\only<3->{
			\begin{ideaBox}
				\textbf{Clave:} los exponentes se multiplican.
			\end{ideaBox}
		}
	\end{frame}
	
	% ----------------------------------------------------------
	\begin{frame}{\iconoDesarrollo\ \ Propiedad 4}
		\centering
		{\Large\bfseries Potencia de un producto}
		
		\vspace{0.35cm}
		
		\[
		(ab)^n=a^n b^n
		\]
		
		\only<2->{
			\begin{exampleblock}{Ejemplo}
				\[
				(2\cdot5)^3=2^3\cdot5^3=8\cdot125=1000
				\]
			\end{exampleblock}
		}
	\end{frame}
	
	% ----------------------------------------------------------
	\begin{frame}{\iconoDesarrollo\ \ Propiedad 5}
		\centering
		{\Large\bfseries Potencia de un cociente}
		
		\vspace{0.35cm}
		
		\[
		\left(\frac{a}{b}\right)^n=\frac{a^n}{b^n},\qquad b\neq0
		\]
		
		\only<2->{
			\begin{exampleblock}{Ejemplo}
				\[
				\left(\frac{2}{3}\right)^3=\frac{2^3}{3^3}=\frac{8}{27}
				\]
			\end{exampleblock}
		}
	\end{frame}
	
	% ----------------------------------------------------------
	\begin{frame}{\iconoDesarrollo\ \ Propiedad 6}
		\centering
		{\Large\bfseries Exponente cero}
		
		\vspace{0.35cm}
		
		\[
		a^0=1,\qquad a\neq0
		\]
		
		\only<2->{
			\begin{exampleblock}{Ejemplo}
				\[
				7^0=1
				\]
			\end{exampleblock}
		}
	\end{frame}
	
	% ----------------------------------------------------------
	\begin{frame}{\iconoDesarrollo\ \ Propiedad 7}
		\centering
		{\Large\bfseries Exponente negativo}
		
		\vspace{0.35cm}
		
		\[
		a^{-n}=\frac{1}{a^n},\qquad a\neq0
		\]
		
		\only<2->{
			\begin{exampleblock}{Ejemplo}
				\[
				2^{-3}=\frac{1}{2^3}=\frac{1}{8}
				\]
			\end{exampleblock}
		}
	\end{frame}
	
	% ----------------------------------------------------------
	\begin{frame}{\iconoDesarrollo\ \ Casos especiales}
		\centering
		\[
		a^1=a
		\qquad\qquad
		1^n=1
		\qquad\qquad
		0^n=0 \text{ para } n>0
		\]
		
		\only<2->{
			\begin{exampleblock}{Ejemplos}
				\[
				9^1=9
				\qquad
				1^{12}=1
				\qquad
				0^5=0
				\]
			\end{exampleblock}
		}
	\end{frame}
	
	% ----------------------------------------------------------
	\SectionFrame{\iconoError}{Errores frecuentes}{Lo que no debemos hacer}
	
	% ----------------------------------------------------------
	\begin{frame}{\iconoError\ \ Error 1}
		\centering
		{\Large\bfseries No sumar exponentes en una suma}
		
		\vspace{0.35cm}
		
		\[
		a^m+a^n \neq a^{m+n}
		\]
		
		\only<2->{
			\begin{exampleblock}{Ejemplo}
				\[
				2^2+2^3=4+8=12
				\]
				\[
				2^{2+3}=2^5=32
				\]
			\end{exampleblock}
		}
	\end{frame}
	
	% ----------------------------------------------------------
	\begin{frame}{\iconoError\ \ Error 2}
		\centering
		{\Large\bfseries No distribuir el exponente en una suma}
		
		\vspace{0.35cm}
		
		\[
		(a+b)^n \neq a^n+b^n
		\]
		
		\only<2->{
			\begin{exampleblock}{Ejemplo}
				\[
				(2+3)^2=5^2=25
				\]
				\[
				2^2+3^2=4+9=13
				\]
			\end{exampleblock}
		}
	\end{frame}
	
	% ----------------------------------------------------------
	\SectionFrame{\iconoEval}{Chequeo rápido}{Verificamos comprensión}
	
	% ----------------------------------------------------------
	\begin{frame}{\iconoEval\ \ Miniverificación}
		\begin{block}{Responde rápido}
			\begin{enumerate}
				\item<1-> \(2^3\cdot2^2=\ ?\)
				\item<2-> \((4^2)^2=\ ?\)
				\item<3-> \(10^0=\ ?\)
				\item<4-> \(3^{-2}=\ ?\)
			\end{enumerate}
		\end{block}
		
		\only<5>{
			\begin{exampleblock}{Respuestas}
				\[
				2^3\cdot2^2=2^5=32
				\qquad
				(4^2)^2=4^4=256
				\]
				\[
				10^0=1
				\qquad
				3^{-2}=\frac{1}{9}
				\]
			\end{exampleblock}
		}
	\end{frame}
	
	% ----------------------------------------------------------
	\SectionFrame{\iconoPractica}{Práctica guiada}{Pasos claros y develados}
	
	% ----------------------------------------------------------
	\begin{frame}{\iconoPractica\ \ Ejercicio guiado 1}
		\begin{block}{Problema}
			Simplificar:
			\[
			2^4\cdot2^3
			\]
		\end{block}
		
		\only<2>{
			\begin{exampleblock}{Paso 1}
				Misma base:
				\[
				2^4\cdot2^3=2^{4+3}
				\]
			\end{exampleblock}
		}
		
		\only<3>{
			\begin{exampleblock}{Paso 2}
				\[
				2^{4+3}=2^7
				\]
			\end{exampleblock}
		}
		
		\only<4>{
			\begin{exampleblock}{Paso 3}
				\[
				2^7=128
				\]
			\end{exampleblock}
		}
	\end{frame}
	
	% ----------------------------------------------------------
	\begin{frame}{\iconoPractica\ \ Ejercicio guiado 2}
		\begin{block}{Problema}
			Resolver:
			\[
			(3^2)^3
			\]
		\end{block}
		
		\only<2>{
			\begin{exampleblock}{Paso 1}
				Potencia de una potencia:
				\[
				(3^2)^3=3^{2\cdot3}
				\]
			\end{exampleblock}
		}
		
		\only<3>{
			\begin{exampleblock}{Paso 2}
				\[
				3^{2\cdot3}=3^6
				\]
			\end{exampleblock}
		}
		
		\only<4>{
			\begin{exampleblock}{Paso 3}
				\[
				3^6=729
				\]
			\end{exampleblock}
		}
	\end{frame}
	
	% ----------------------------------------------------------
	\SectionFrame{\iconoAgro}{Aplicaciones agronómicas}{Conexión profesional}
	
	% ----------------------------------------------------------
	\begin{frame}{\iconoAgro\ \ Caso 1: semillas por superficie}
		\begin{block}{Situación}
			Se usan \(2^5\) semillas por metro cuadrado en una parcela de \(2^3\) metros cuadrados.
		\end{block}
		
		\only<2>{
			\begin{exampleblock}{Paso 1}
				\[
				2^5\cdot2^3
				\]
			\end{exampleblock}
		}
		
		\only<3>{
			\begin{exampleblock}{Paso 2}
				\[
				2^5\cdot2^3=2^8
				\]
			\end{exampleblock}
		}
		
		\only<4>{
			\begin{exampleblock}{Resultado}
				\[
				2^8=256
				\]
				Se utilizan \textbf{256 semillas}.
			\end{exampleblock}
		}
	\end{frame}
	
	% ----------------------------------------------------------
	\begin{frame}{\iconoAgro\ \ Caso 2: crecimiento bacteriano}
		\begin{block}{Situación}
			Una colonia bacteriana parte con \(2^3\) bacterias y se duplica durante 4 horas.
		\end{block}
		
		\only<2>{
			\begin{exampleblock}{Modelo}
				\[
				2^3\cdot2^4
				\]
			\end{exampleblock}
		}
		
		\only<3>{
			\begin{exampleblock}{Cálculo}
				\[
				2^3\cdot2^4=2^7
				\]
			\end{exampleblock}
		}
		
		\only<4>{
			\begin{exampleblock}{Resultado}
				\[
				2^7=128
				\]
				Hay \textbf{128 bacterias}.
			\end{exampleblock}
		}
	\end{frame}
	
	% ----------------------------------------------------------
	\begin{frame}{\iconoAgro\ \ Caso 3: fertilización}
		\begin{block}{Situación}
			Se aplican \(5\times10^2\) gramos por hectárea en una superficie de \(10^3\) hectáreas.
		\end{block}
		
		\only<2>{
			\begin{exampleblock}{Modelo}
				\[
				(5\times10^2)(10^3)
				\]
			\end{exampleblock}
		}
		
		\only<3>{
			\begin{exampleblock}{Cálculo}
				\[
				(5\times10^2)(10^3)=5\times10^5
				\]
			\end{exampleblock}
		}
		
		\only<4>{
			\begin{exampleblock}{Resultado}
				\[
				5\times10^5=500000
				\]
				Se requieren \textbf{500000 gramos}.
			\end{exampleblock}
		}
	\end{frame}
	
	% ----------------------------------------------------------
	\SectionFrame{\iconoPractica}{Práctica autónoma}{Consolidación}
	
	% ----------------------------------------------------------
	\begin{frame}{\iconoPractica\ \ Ejercicios}
		\begin{enumerate}
			\item \(2^5\cdot2^2\)
			\item \(\dfrac{3^7}{3^4}\)
			\item \((4^2)^3\)
			\item \((2\cdot3)^2\)
			\item \(\left(\dfrac{3}{5}\right)^2\)
			\item \(8^0\)
			\item \(2^{-4}\)
			\item En un cultivo hay \(3^4\) plantas por sector y \(3^2\) sectores. ¿Cuántas plantas hay en total?
		\end{enumerate}
	\end{frame}
	
	% ----------------------------------------------------------
	\begin{frame}{\iconoPractica\ \ Respuestas}
		\begin{multicols}{2}
			\begin{enumerate}
				\item \(2^5\cdot2^2=2^7=128\)
				\item \(\dfrac{3^7}{3^4}=3^3=27\)
				\item \((4^2)^3=4^6=4096\)
				\item \((2\cdot3)^2=36\)
				\item \(\left(\dfrac{3}{5}\right)^2=\dfrac{9}{25}\)
				\item \(8^0=1\)
				\item \(2^{-4}=\dfrac{1}{16}\)
				\item \(3^4\cdot3^2=3^6=729\)
			\end{enumerate}
		\end{multicols}
	\end{frame}
	
	% ----------------------------------------------------------
	\SectionFrame{\iconoCierre}{Cierre}{Síntesis y proyección}
	
	% ----------------------------------------------------------
	\begin{frame}{\iconoCierre\ \ Ideas clave}
		\begin{block}{Síntesis}
			\begin{itemize}
				\item<1-> Una potencia representa una multiplicación repetida.
				\item<2-> La base se repite y el exponente cuenta repeticiones.
				\item<3-> Las propiedades simplifican cálculos.
				\item<4-> Las potencias modelan crecimiento y escalas.
				\item<5-> En Agronomía ayudan a interpretar datos reales.
			\end{itemize}
		\end{block}
	\end{frame}
	
	% ----------------------------------------------------------
	\begin{frame}{\iconoTicket\ \ Ticket de salida}
		\begin{exampleblock}{Responde antes de salir}
			\begin{enumerate}
				\item Simplifica:
				\[
				2^3\cdot2^4
				\]
				\item Explica con tus palabras qué significa:
				\[
				10^{-2}
				\]
				\item Nombra una situación de Agronomía donde usarías potencias.
			\end{enumerate}
		\end{exampleblock}
	\end{frame}
	
	% ----------------------------------------------------------
	\begin{frame}{\iconoPuente\ \ Puente hacia la próxima clase}
		\begin{block}{Conexión con notación científica}
			Las potencias de 10 permiten representar números:
			\begin{itemize}
				\item muy grandes, como \(450000=4{,}5\times10^5\),
				\item muy pequeños, como \(0{,}00045=4{,}5\times10^{-4}\).
			\end{itemize}
		\end{block}
		
		\pp
		
		\begin{ideaBox}
			La próxima clase puede continuar naturalmente con \destacar{notación científica}, multiplicación y división usando potencias de 10.
		\end{ideaBox}
	\end{frame}
	
	% ----------------------------------------------------------
	\begin{frame}{\iconoCierre\ \ Cierre final}
		\begin{center}
			{\Huge\color{USTTeal}\faSeedling}
			
			\vspace{0.35cm}
			
			{\Large\bfseries\color{USTBlueDark}Potencias para comprender y modelar}
			
			\vspace{0.35cm}
			
			{\normalsize\textcolor{USTGray}{La matemática básica no solo permite calcular: permite leer mejor el mundo técnico y profesional.}}
		\end{center}
	\end{frame}
	
\end{document}